\label{sec:conclusion}
We studied \emph{posture-requested} personalized text-to-image generation---a
single reference identity, an explicit textual posture request, and
accountability for realizing both---a regime where adapter-based personalizers,
bound by their backbone's pose prior, either fail to rotate or lose identity. We first verified causally that FLUX.2-klein's frozen VAE and joint attention
actually carry and use identity content---not merely occupy a sequence
slot---establishing its in-context conditioning as a valid, if unintentional,
identity substrate. On
FLUX.2-klein we combined an in-context sam3d-skeleton control with reference
duplication to attain strong rotation while recovering identity. We then gave a
mechanistic account grounded in the backbone's own token geometry---references
enter the joint attention as positionally-ordered key/value blocks, so identity
and pose compete for one softmax-normalized attention budget---and confirmed it
causally by attention surgery, turning a collection of training-free heuristics
(duplication, ordering) into a single principle. That principle yielded a
$\sim$200-parameter learned attention-budget allocator that \emph{refines}
identity further---the first trained component to improve over the strongest
training-free configuration at no pose cost. The second half of the contribution
is measurement: our \emph{pose-rectified performance metrics}---an explainable
pose check that verifies whether the requested posture was realized, a
pose-adapted identity read-out (PAEs), and a pose-aware identity that credits the
face cosine only at the posture actually produced---both expose the profile-view
identity loss central to the task and prevent frontal collapse from masquerading
as identity preservation. Future work includes scaling the learned allocator to
many more identities via a self-distilled data engine, testing it on a
non-distilled backbone, and closing the chin-down pitch axis with pose-targeted
data or 3D-aware control.
